% Drawing 154/03-16 -- folder 154, batch 03, archive page 55 (PDF page 56).
% Transcription: figure : un cercle au crayon, en gros tirets, marqué « donc » à gauche
% Kind: dashed pencil circle. Method 2 (schematic redraw from the facsimile). Status: ok.
% Notes: Five thick pencil dashes on a circle (a long one at the top, two short ones at the sides, two at the bottom); the word « donc » at its left; the formula written inside the circle belongs to the text. Redrawn 2026-11-12 (issue 128): to check again.
% Coordinates: 1 unit = 1 mm = 1/100 of the crop width (bbox in page fractions: [0.6, 0.255, 0.8, 0.365]).
% Redrawn by Claude (claude-fable-5-2), issue #128, 2026-11-12; not checked by a human; not compiled (no TeX in the session).
% Rights: reproduces part of the work -- behind RENDER_TRANSCRIPTS, never published by itself.
\documentclass[tikz,border=3pt]{standalone}
\usepackage{amsmath,amssymb}
\usetikzlibrary{patterns,arrows.meta}
\tikzset{pencil/.style={black!45}, struck/.style={black!35}, mpt/.style={circle,draw,inner sep=1.3pt}, pt/.style={circle,fill,inner sep=0.9pt}, opt/.style={circle,draw,fill=white,inner sep=1.1pt}}
\begin{document}
\begin{tikzpicture}[x=1mm,y=1mm,line cap=round,line join=round,font=\scriptsize]
\draw[pencil,very thick] (64.3,-25.6) arc[start angle=39.8, end angle=131.3, x radius=26.1, y radius=26.1];
\draw[pencil,very thick] (18.2,-40.4) arc[start angle=175.9, end angle=192.5, x radius=26.1, y radius=26.1];
\draw[pencil,very thick] (69.7,-47.8) arc[start angle=347.9, end angle=362, x radius=26.1, y radius=26.1];
\draw[pencil,very thick] (26.6,-61.6) arc[start angle=227.6, end angle=249.2, x radius=26.1, y radius=26.1];
\draw[pencil,very thick] (53,-66.9) arc[start angle=289.7, end angle=312.3, x radius=26.1, y radius=26.1];
\node[inner sep=1pt] at (11.8,-41.4) {\small donc};
\end{tikzpicture}
\end{document}
